\section{Uniform random-matrix estimates}\label{sec:matrices}

The field analysis uses random-matrix estimates in two forms.  Finite
field recursions produce products whose diagonal coefficients may be
chosen after the disorder is observed.  Lemma~\ref{lem:words} controls
such products of bounded length, subject to a separate positivity
margin whenever an inverse occurs.  We also obtain fixed-vector
bilinear estimates and a simultaneous determinant bound for the
stability argument in Section~\ref{sec:stability}.  Their expectation
calculations use the same inverse identity.  We begin with fixed
deterministic coefficients and then use concentration and finite meshes
to obtain the required simultaneous estimates.

Throughout this section, $j\in(0,1)$ is fixed.  A centered GOE matrix
$W$ has independent entries on and above the diagonal, with
\[
 \E W_{ab}^{2}=j/n\quad(a<b),\qquad
 \E W_{aa}^{2}=2j/n.
\]
We couple $J$ to $W$ by putting $J=W-D_{\diag W}$.  In particular,
the off-diagonal law of $J$ is the law used in the theorem.  Fix
$\eta>0$ so that, with $a_+=1+\eta$,
\begin{equation}\label{eq:matrix-subcritical}
 j a_+^2<1.
\end{equation}
For a nonnegative diagonal matrix $A\le a_+I$, put
\begin{equation}\label{eq:K}
 \begin{gathered}
 \chi_A=\frac1n\Tr A,\qquad B_A=j\chi_A,\qquad
 K_A(M)=K_A(1;M),\\
 K_A(z;M)=(I-zAM+z^2B_AA)^{-1}.
 \end{gathered}
\end{equation}
The associated symmetric matrix is
\[
 S_A(z;M)=I+z^2B_AA-zA^{1/2}MA^{1/2}.
\]
The symmetric matrix lets us control the inverse without imposing a
positive lower bound on the individual diagonal entries of $A$.
Whenever $S_A(z;M)$ is invertible, the following identities hold:
\begin{align}
 K_A(z;M)
 &=I+A^{1/2}S_A(z;M)^{-1}A^{1/2}(zM-z^2B_AI),
 \label{eq:K-symmetric-representation}\\
 K_A(z;M)A&=A^{1/2}S_A(z;M)^{-1}A^{1/2}.
 \label{eq:KA-symmetric-representation}
\end{align}
To verify the first identity, apply
$(I-UV)^{-1}=I+U(I-VU)^{-1}V$.  For the second, multiply the two
sides using $S_A=I-A^{1/2}(zM-z^2B_AI)A^{1/2}$.  These formulas
will also let us truncate the inverse while retaining bounded
matrix factors.

\subsection{Words and their predictions}

A \emph{word} is an ordered product of diagonal matrices, copies of a
symmetric matrix $M$, and at most one factor $K_A(M)$.  Its length is
the number of factors.  Consecutive diagonal factors may always be
combined, and an absent diagonal factor is interpreted as $I$.
To each word $Q$ we associate a diagonal matrix $\mathsf P(Q)$ by
the contraction rule below.  The estimates will compare the diagonal
of $Q$ with this explicitly defined matrix.

Begin with a word without $K_A$.  Consider every complete
noncrossing pairing of its occurrences of $M$, in their linear
order.  For each pairing, remove an innermost pair by the rule
\begin{equation}\label{eq:noncrossing-contraction}
 MDM\ \longmapsto\ j\frac{\Tr D}{n}I,
\end{equation}
where $D$ is diagonal after previous removals.  Combine adjacent
diagonals and continue.  Sum the resulting diagonal matrices over
all complete noncrossing pairings.  If there is no complete pairing,
the prediction is zero; if there is no occurrence of $M$, it is the
word itself.  The result for a fixed pairing does not depend on the
order of removal: two available innermost pairs are disjoint, and
their removals commute.
For words without an inverse factor, this is the operator-valued
semicircular moment rule with covariance map
$D\mapsto jn^{-1}\Tr(D)I$; see Speicher~\cite{Speicher1998}.

For example, four occurrences of $M$ have two complete noncrossing
pairings, giving
\[
 \begin{split}
 \mathsf P(MD_1MD_2MD_3M)
 ={}&j^2\frac{\Tr D_1}{n}\frac{\Tr D_3}{n}D_2\\
 &+j^2\frac{\Tr D_2}{n}\frac{\Tr(D_1D_3)}n I.
 \end{split}
\]
The first term pairs the first two and last two occurrences; the second
pairs the outside occurrences around the middle pair.

For a word containing $K_A$, replace that factor by the formal series
\begin{equation}\label{eq:K-formal-series}
 K_A(z;M)=\sum_{r\ge0}(zAM-z^2B_AA)^r,
\end{equation}
and apply the preceding rule coefficient by coefficient.  Ordinary
occurrences of $M$ outside this factor are left unscaled.  We prove
below that the resulting series $\mathsf P(Q;z)$ is a polynomial;
we define $\mathsf P(Q)=\mathsf P(Q;1)$.

The corresponding trace prediction is
$\mathsf p(Q)=n^{-1}\Tr\mathsf P(Q)$.  To see its cyclic invariance,
draw an ordinary trace word on a circle, with a chord for each
pair in a noncrossing pairing.  Its contribution is $j$ for each chord,
multiplied by, for each face, the average of the product of the
diagonal matrices on that face's boundary arcs.  Removing an
innermost chord gives exactly \eqref{eq:noncrossing-contraction}, and
the average over the last face is the final normalized trace.  This
description is unchanged by rotating the circle.  For words with
$K_A(z;M)$, apply this argument coefficient by coefficient in
\eqref{eq:K-formal-series}.

\begin{lemma}[Simultaneous word estimates]\label{lem:words}
Fix $L\in\mathbb N$, $M_0<\infty$, and $c>0$.  There are constants
$C,c_1>0$ and $n_0$, depending only on $j,a_+,L,M_0,c$, such that the
following event has probability at least $1-Ce^{-c_1n}$ for
$n\ge n_0$.  Simultaneously for all words $Q$ of length at most $L$,
all their real diagonal factors of norm at most $M_0$, and all
$0\le A\le a_+I$, one has
\begin{equation}\label{eq:word-estimates}
 \norm{Q}\le C,\qquad
 \left(\sum_{a\ne b}|Q_{ab}|^4\right)^{1/4}\le C,\qquad
 \norm{\diag(Q-\mathsf P(Q))}_2\le C.
\end{equation}
For a word containing $K_A(J)$, the assertion is restricted to
parameters satisfying $S_A(1;J)\succeq cI$.  The same conclusion
holds with $W$ in place of $J$, with the corresponding stability
restriction.  Thus all diagonal parameters in this assertion may
be chosen after the random matrix has been observed.

If a word has $r$ ordinary occurrences of $M$ outside its one
$K_A(z;M)$ factor, then $\mathsf P(Q;z)$ has degree at most $r$.
\end{lemma}

We prove the simultaneous estimate in three steps.  For a fixed
deterministic diagonal, Gaussian comparison gives stability along
$0\le z\le1$.  We then calculate the expected diagonal of a
truncated word.  Finally, concentration makes the word estimates
simultaneous in the diagonal parameters.  The first step supplies
stability for each deterministic choice; it does not remove the
separate stability restriction for adaptive $A$ in
Lemma~\ref{lem:words}.

\begin{lemma}[Stability for a deterministic diagonal]
\label{lem:matrix-path-stability}
There is $c_0>0$, depending only on $j,a_+$, with the following
property.  For every fixed deterministic $0\le A\le a_+I$, and every
fixed $R>2\sqrt j$, there are constants $C,c_2>0$ and $n_0$,
uniform in $A$, such that
\begin{equation}\label{eq:deterministic-path-stability}
 \Prob\left\{\norm W\le R\ \hbox{and}\quad
 S_A(z;W)\succeq c_0I\text{ for all }0\le z\le1\right\}
 \ge1-Ce^{-c_2n}\qquad(n\ge n_0).
\end{equation}
In addition, for every fixed $\varepsilon>0$,
$\Prob\{\norm W>2\sqrt j+\varepsilon\}\le
C_\varepsilon e^{-c_\varepsilon n}$ for sufficiently large $n$.
\end{lemma}

\begin{proof}
We first bound the expected quadratic process that defines the
smallest eigenvalue of $S_A$.  Concentration will give a margin at
each fixed $z$, and a finite mesh will extend it along the path.

We use two standard Gaussian facts.  A real-valued
$L$-Lipschitz function of a standard Gaussian vector satisfies
$\Prob\{|F-\E F|>u\}\le2\exp(-u^2/(2L^2))$.
Also, if two centered Gaussian processes have ordered squared
increments, the process with larger squared increments has the
larger expected supremum (Sudakov--Fernique). The latter comparison remains valid
after the same deterministic function is added to both processes;
see the deterministic-shift comparison in Vitale~\cite{Vitale2000}.
Both facts apply by finite nets and continuity to the compact
index sets below.  The linear map from the independent standard
Gaussian coordinates to $W$, in Frobenius norm, has norm
$\sqrt{2j/n}$.

For the comparison, write $w=A^{1/2}p$ for $\norm p=1$ and consider
\[
 X_p=w^TWw,\qquad
 Y_p=2\sqrt{j/n}\,\norm w\,g^Tw,
\]
where $g$ is standard Gaussian in $\R^n$.  If $w'=A^{1/2}p'$, direct
calculation gives
\begin{align*}
 &\E(Y_p-Y_{p'})^2-\E(X_p-X_{p'})^2\\
 &\qquad=\frac{2j}{n}
 \left[(\norm w^2-\norm{w'}^2)^2
       +2(\norm w\norm{w'}-w^Tw')^2\right]\ge0.
\end{align*}
Consequently, for each $z\in[0,1]$,
\[
 \E\sup_{\norm p=1}\{zX_p-z^2j\chi_A\norm w^2\}
 \le \E\sup_{\norm p=1}\{zY_p-z^2j\chi_A\norm w^2\}.
\]
Set $T_g=\sqrt{j/n}\norm{A^{1/2}g}$.  Since
$g^Tw=p^TA^{1/2}g\le\norm{A^{1/2}g}$, the last supremum is at most
\[
 \sup_{0\le v\le\sqrt{a_+}}
       (2zT_gv-z^2j\chi_Av^2).
\]
The remaining supremum depends on the Gaussian vector only through
$T_g$.  The variable $T_g^2$ has mean $j\chi_A$ and variance
$2j^2\Tr(A^2)/n^2\le C/n$.  The inequality
$|\sqrt x-\sqrt y|\le\sqrt{|x-y|}$ therefore shows, uniformly in
$A$, that $\E|T_g-\sqrt{j\chi_A}|=o(1)$.  On replacing $T_g$ by
$\sqrt{j\chi_A}$ in the preceding display, the expression becomes
$2d-d^2$, with
$0\le d=z\sqrt{j\chi_A}v\le\sqrt j\,a_+<1$.
Thus, with $\delta=(1-\sqrt j\,a_+)^2>0$,
\begin{equation}\label{eq:quadratic-process-margin}
 \E\sup_{\norm p=1}
 \{zp^TA^{1/2}WA^{1/2}p-z^2j\chi_Ap^TAp\}
 \le1-\delta+o(1),
\end{equation}
uniformly in $A,z$.
The supremum in \eqref{eq:quadratic-process-margin} is
$a_+$-Lipschitz in $W$ for Frobenius norm.  Gaussian concentration
gives an exponentially small probability that it exceeds
$1-\delta/2$, uniformly for each fixed $A,z$.

To pass from a fixed $z$ to the entire interval, we also need a
bound on $\norm W$.  Taking $A=I$ in the comparison without
deterministic offsets gives
\[
 \E\lambda_{\max}(W)\le2\sqrt{j/n}\E\norm g\le2\sqrt j.
\]
The same holds for $-W$.  Concentration of these two eigenvalues
proves the stated norm tail.  On $\norm W\le R$, the function of
$z$ inside the supremum in \eqref{eq:quadratic-process-margin} is
uniformly Lipschitz, with constant at most $a_+R+2ja_+^2$.
A sufficiently fine constant mesh of $[0,1]$, followed by a union
bound, therefore gives \eqref{eq:deterministic-path-stability}, for
example with $c_0=\delta/4$.
\end{proof}

\subsection{Truncation and the loop recursion}

The inverse in a word need not exist for every realization of the
matrix.  We therefore replace each word by a globally defined
bounded function of $W$, agreeing with the original word whenever
the required norm and stability conditions hold.  This permits
concentration and integration by parts before those conditions
are imposed.

Fix a sufficiently large constant $R$.  Let $\Pi_R$ be metric
projection, in the Hilbert space of real symmetric matrices with
Frobenius norm, onto $\{M:\norm M\le R\}$.  This is a nonexpansive
map and is equivariant under orthogonal conjugation.  Put
$\widehat W=\Pi_R(W)$.  Choose a smooth compactly supported real
function $f$ which equals $x^{-1}$ on the interval
\[
 \left[\tfrac12\min(c,c_0),\,2+c+c_0+ja_+^2+a_+R\right].
\]
This interval contains, with slack, every stable spectrum needed
both for deterministic-parameter estimates and for the restriction
$S_A(1;M)\succeq cI$ in Lemma~\ref{lem:words}.  Replace $K_A(z;W)$
by
\begin{equation}\label{eq:truncated-K}
 \widehat K_A(z)
 =I+A^{1/2}f(S_A(z;\widehat W))A^{1/2}
       (z\widehat W-z^2B_AI).
\end{equation}
Replace every ordinary $W$ in a word by $\widehat W$ as well; denote
the resulting word by $\widehat Q_\theta(z)$, where $\theta$ lists
all its diagonal parameters.  These truncated factors equal the
original factors on the stability events where they will be used.

For each fixed length and fixed bounds on the diagonal parameters,
there is a constant $C$ such that
\begin{align}
 \norm{\widehat Q_\theta(z)}&\le C,
 &\norm{\widehat Q_\theta(z;W)-\widehat Q_\theta(z;W')}_{\Frob}
 &\le C\norm{W-W'}_{\Frob},
 \label{eq:truncated-word-lipschitz}\\
 \Lip_{\Frob}\bigl(\widehat Q_\theta(z;\cdot)
                   -\widehat Q_{\theta'}(z;\cdot)\bigr)
 &\le C\sqrt\Delta
 &&\text{if }\norm{\theta-\theta'}_\infty\le\Delta\le1.
 \label{eq:truncated-word-increment-lipschitz}
\end{align}
The constants are uniform for $z\in[0,1]$ and for all matrix sizes.
Here the Lipschitz constants concern maps into matrices with
Frobenius norm.

We verify both estimates with constants independent of dimension.
On the operator-norm ball, all factors are bounded, and their
differentials in a perturbation $E$ are bounded in Frobenius norm
by $C\norm E_{\Frob}$.  The only functional-calculus factor is
handled by writing $f$ as a Fourier integral and using
\[
 \dd(e^{itS})[E]
 =it\int_0^1 e^{it(1-u)S}E e^{ituS}\,\dd u.
\]
The Fourier transform of $f$ decays faster than any power, so its
first two weighted absolute moments are finite.  This formula and
the corresponding difference formula give
\[
 \norm{\dd f(S)[E]}_{\Frob}\le C\norm E_{\Frob},\qquad
 \norm{(\dd f(S)-\dd f(S'))[E]}_{\Frob}
 \le C\norm{S-S'}\norm E_{\Frob}.
\]
Changing the diagonal parameters by at most $\Delta$ changes their
square roots in norm by at most $\sqrt\Delta$.  Consequently the
factors, and their matrix differentials, change by respectively
$C\sqrt\Delta$ and $C\sqrt\Delta\norm E_{\Frob}$.  The product rule
proves the same differential assertions for a word.  Integrating
along the line segment between two matrices in the convex
operator-norm ball proves \eqref{eq:truncated-word-lipschitz} and
\eqref{eq:truncated-word-increment-lipschitz} there.  Composition with
the nonexpansive projection $\Pi_R$ proves them everywhere.

These Lipschitz estimates first give the moment bounds needed for
the expectation calculation.  For a fixed parameter tuple,
\begin{equation}\label{eq:word-entry-and-trace-variance}
 \Var(\widehat Q_{ab})\le C/n,\qquad
 \Var\left(\frac1n\Tr\widehat Q\right)\le C/n^2.
\end{equation}
The normalized trace gains a factor $n^{-1/2}$ in its Frobenius
Lipschitz constant.  The same estimates for a difference of two
parameter tuples gain a factor $\Delta$ on the right-hand sides.
Fourth centered entry moments are bounded by $C/n^2$, or
$C\Delta^2/n^2$ for such a difference, by integrating the Gaussian
concentration tail.

\begin{lemma}[Fixed-parameter diagonal expectations]
\label{lem:word-bias}
For the truncated words just defined, uniformly in deterministic
diagonal parameters in the indicated bounded ranges, in
$z\in[0,1]$, and in the diagonal index $i$,
\begin{equation}\label{eq:word-bias}
 \E\widehat Q_\theta(z)_{ii}
 =\mathsf P(Q_\theta;z)_{ii}+O(n^{-1}).
\end{equation}
For words without $K_A$, the prediction has no $z$ dependence.
The formal prediction for a word with $r$ ordinary $W$ letters is a
polynomial of degree at most $r$.
\end{lemma}

\begin{proof}
The calculation has two parts.  We first justify integration by
parts for the truncated words and compute the inverse factor
alone.  We then obtain a recursion in the number of ordinary
$W$ factors; the same recursion determines the formal
prediction exactly and the expectation up to $O(n^{-1})$.

All constants in this proof depend only on the fixed length and
parameter bounds.  Lemma~\ref{lem:matrix-path-stability} shows that,
for every fixed deterministic parameter tuple, the truncations
agree with the actual matrices for all $z\in[0,1]$, except on an
event of probability $Ce^{-c_2n}$.  Gaussian integration by parts
can be applied to the globally Lipschitz truncated factors, using
weak derivatives.  More explicitly, for a scalar Lipschitz
function $H$,
\begin{equation}\label{eq:GOE-integration-by-parts}
 \E[W_{ab}H(W)]
 =\frac jn\E\bigl[\dd H(W)[E_{ab}+E_{ba}]\bigr].
\end{equation}
When $a=b$, the elementary direction on the right is doubled,
as required by the diagonal variance $2j/n$.

For \eqref{eq:GOE-integration-by-parts}, remove the projection
from the distinguished ordinary $W$ factor.  The resulting
expectation error is exponentially small, since the other factors
are bounded and Gaussian norm tails control the removed factor.
On the event of agreement, the differential of the inverse is
\begin{equation}\label{eq:inverse-differential}
 \dd K_A(z;W)[E]=zK_A(z;W)AEK_A(z;W).
\end{equation}
We may use the expressions on the right, with truncated factors,
on the whole probability space.  The error is exponentially small
even after the index sums: all factors and their true truncated
differentials are bounded, there are at most a fixed polynomial
number of indices, and the disagreement event has exponentially
small probability.  More precisely, the elementary direction
$E_{ab}+E_{ba}$ has Frobenius norm at most $2$.  Both the true weak
differential and its formal replacement are therefore bounded by
$C$ in each entry; their difference vanishes on the agreement event.
Each such differential error has expected absolute value at most
$Ce^{-c_2n}$.  Expanding the diagonal entry of a length-$L$ word and
applying the product rule uses at most
$(L-1)n^{L-1}$ scalar contributions before any index identifications;
the factor $j/n$ in integration by parts only improves this bound.
Thus their total error is at most a fixed polynomial in $n$ times
$e^{-c_2n}$, hence at most $C'e^{-c'n}$.  Subsequent words have
bounded lengths depending only on the original length.  We
therefore suppress hats during the loop calculation, with this
convention understood.

We now compute the base case, in which the only nondiagonal
factor is $K=K_A(z;W)$.  Write
$a_i=A_{ii}$, $k_i(z)=\E K_{ii}$, and
$r(z)=n^{-1}\sum_i a_i k_i(z)$.  Differentiating $K_{bi}$ in
\eqref{eq:GOE-integration-by-parts} gives
\begin{align*}
 \E(WK)_{ii}
 &=\frac{zj}{n}\E\left[(\Tr KA)K_{ii}\right]
   +\frac{zj}{n}\sum_b\E[(KA)_{bi}K_{bi}]+O(e^{-c'n})\\
 &=zjr(z)k_i(z)+O(n^{-1}).
\end{align*}
For the second sum, the absolute value of the sum over $b$ is at
most $\norm{KAe_i}\norm{Ke_i}\le C$.  The factorization in the
first term follows from \eqref{eq:word-entry-and-trace-variance};
its covariance is at most $C n^{-3/2}$, which is more than enough.
The inverse identity $K=I+zAWK-z^2B_AAK$ now yields
\begin{equation}\label{eq:K-loop-equations}
 \E(WK)_{ii}=zjr k_i+O(n^{-1}),\qquad
 k_i=1+z^2ja_i(r-\chi_A)k_i+O(n^{-1}).
\end{equation}

The loop equations must select the solution near $k_i=1$.
We establish this by continuation from $z=0$, where the inverse
is the identity.  Put
$h(z)=r(z)-\chi_A$.  The truncated $k_i$ are uniformly bounded and
continuous in $z$.  If $|h(z)|\le\delta_1$, the second equation in
\eqref{eq:K-loop-equations} gives
$\max_i|k_i(z)-1|\le C\delta_1+C/n$.  Choose $\delta_1>0$ small,
using \eqref{eq:matrix-subcritical}, so that throughout this region
and for large $n$,
\[
 \left|z^2j\frac1n\sum_i a_i^2 k_i(z)\right|
 \le1-\delta_2
\]
for some fixed $\delta_2>0$.  Averaging the second equation in
\eqref{eq:K-loop-equations} with weights $a_i$ gives
\[
 h(z)=z^2j\left(\frac1n\sum_i a_i^2k_i(z)\right)h(z)+O(n^{-1}),
\]
and hence $|h(z)|\le C/n$ while $|h(z)|\le\delta_1$.  Since
$h(0)=0$, continuity rules out a first exit from that region for
large $n$.  It follows uniformly in $z$ that
\begin{equation}\label{eq:K-and-WK-bias}
 \E K_{ii}=1+O(n^{-1}),\qquad
 \E(WK)_{ii}=zB_A+O(n^{-1}).
\end{equation}

To complete the base case, we verify the formal identity
$\mathsf P(K_A(z;W))=I$ independently of the expectation calculation.
The cancellation takes place at each coefficient of the series.
In the expansion
\eqref{eq:K-formal-series}, a block $-z^2B_AA$ may be regarded as
the negative contraction of two adjacent blocks $zAW,zAW$, since
their contraction is $z^2j\chi_AA=z^2B_AA$.
Fix a completely noncrossing-paired chain $(AW)^{2m}$ with $m>0$.
Let $\mathcal A$ be the nonempty set of its pairs which join two
originally adjacent occurrences of $W$.  It is nonempty because
every nonempty noncrossing pairing of a linear order has an
adjacent pair.  All pairs in $\mathcal A$ are disjoint.  The terms
of \eqref{eq:K-formal-series} which give this paired chain are in
bijection with subsets $T\subseteq\mathcal A$: contract the pairs
in $T$ into the negative deterministic blocks, and pair the
remaining occurrences normally.  The magnitude of the final
diagonal contribution is unchanged, and its sign is $(-1)^{|T|}$.
Thus the total multiplier is
$\sum_{T\subseteq\mathcal A}(-1)^{|T|}=0$.  This argument is finite
at each power of $z$ and proves the asserted formal identity.
Multiplication by surrounding diagonal matrices proves both
claims of the lemma for words with no ordinary $W$ letters.

For the induction step, select an ordinary occurrence of $W$ in
a general word and apply integration by parts at that occurrence.
We describe the resulting finite recursion directly, without
requiring convergence of the formal series.  In each product-rule
term, mark the partner of the selected occurrence: for a derivative
falling on an ordinary factor, this is the $W$ being differentiated.
If the derivative falls on $K_A$, replace that factor by
$zK_AA\,\boxed W\,K_A$ and take the boxed occurrence as the
partner.  In either case, list the two marked positions in their
linear order, so the marked product reads
\[
 U\,\boxed W\,V\,\boxed W\,Q.
\]
The two orientations in \eqref{eq:GOE-integration-by-parts}
contribute respectively
\begin{equation}\label{eq:word-loop-contractions}
 \frac jn\E\bigl[(UQ)_{ii}\Tr V\bigr],\qquad
 \frac jn\E(UV^TQ)_{ii}.
\end{equation}
There is an additional factor $z$ when the derivative fell on
$K_A$.  The second expression in
\eqref{eq:word-loop-contractions} is $O(n^{-1})$ by operator-norm
bounds, including when it contains two copies of $K_A$.
The first expression is
\begin{equation}\label{eq:word-loop-leading}
 j\left(\frac1n\E\Tr V\right)\E(UQ)_{ii}+O(n^{-1}),
\end{equation}
again by \eqref{eq:word-entry-and-trace-variance}.  The index sums
explain the two terms.  The opposite orientation
identifies the two interior indices and gives $\Tr V$; the equal
orientation reverses the interior path and gives $V^T$.

Consequently $\E\widehat Q_{ii}$ is the finite sum of the leading
products \eqref{eq:word-loop-leading} over all ordinary partners,
plus $z$ times the corresponding leading product for the one
$K_A$ factor if present, with an $O(n^{-1})$ remainder.  Each word
$V$ and $UQ$ in this recursion has fewer ordinary $W$ letters than
the original word.  Each contains at most one $K_A$: if the
derivative fell on $K_A$, the splitting places one of its two
copies in each word.  Extra diagonal factors $A$ cause no problem,
since at most one is added per decrease of the induction index.

The decrease in the number of ordinary letters holds whether the
selected occurrence lies before or after the inverse.  The two
possible orders can be checked explicitly.  If the original word is
$P K_A R W T$, the contraction with $K_A$ contributes
\[
 zj\,\E\left[(P K_A A T)_{ii}\frac{\Tr(K_A R)}n\right].
\]
If it is $P W R K_A T$, that contribution is
\[
 zj\,\E\left[(P K_A T)_{ii}\frac{\Tr(R K_A A)}n\right].
\]
Here $P,R,T$ contain no inverse factor.  The displayed outer and
inner words each have one inverse factor, even though the original
distinguished occurrence may lie anywhere among the ordinary letters.

It remains to compare this expectation recursion with the formal
prediction.  Sort the noncrossing pairings by the partner of the
selected ordinary occurrence.  The inside and outside may then be paired
independently; their predictions are respectively
$\mathsf p(V)$ and $\mathsf P(UQ)$.  If the partner occurs inside
the series for $K_A(z;W)$, marking that occurrence in the formal
series is exactly the formal differentiation identity
$\dd K_A[E]=zK_AAEK_A$.  Thus the entire sum over positions in that
series is the single $z$-weighted product just described.

We can now induct on the number $r$ of ordinary $W$ letters.
An ordinary partner leaves a total of $r-2$ such letters in the
two smaller words; a partner in $K_A$ leaves $r-1$.  The base
prediction is constant.  The recursion therefore proves that
the formal prediction is a polynomial of degree at most $r$.
It also proves \eqref{eq:word-bias}: use the induction hypothesis
for the two expectations in each leading product, and average
the diagonal induction hypothesis to handle its normalized
trace.  There are only finitely many products and all factors
are uniformly bounded.  The same induction, with no inverse
factor, proves the statement for ordinary words.
\end{proof}

\subsection{Making the word estimates simultaneous}

\begin{proof}[Proof of Lemma~\ref{lem:words}]
We first make the centered estimates uniform over the diagonal
parameters, then add the expectation estimate from
Lemma~\ref{lem:word-bias}.  Fix a letter pattern and let $\theta$
range over its parameter box, which has at most $C_Ln$ real
coordinates.  Work initially with the truncated words at $z=1$.  Conjugation by any
diagonal sign matrix preserves the law of $W$ and conjugates
$\widehat Q_\theta$ by that same sign matrix.  The diagonal
parameters commute with the conjugation, and both the projection
and the functional calculus respect it.  Thus
$\E(\widehat Q_\theta)_{ab}=0$ whenever $a\ne b$.

Define the two matrix seminorms
\[
 N_2(M)=\norm{\diag M}_2,\qquad
 N_4(M)=\left(\sum_{a\ne b}|M_{ab}|^4\right)^{1/4}.
\]
Both are $1$-Lipschitz with respect to Frobenius norm.
The entry moment bounds following
\eqref{eq:word-entry-and-trace-variance} give
\begin{equation}\label{eq:centered-word-norm-expectations}
 \E N_k(\widehat Q_\theta-\E\widehat Q_\theta)\le C
 \quad(k=2,4).
\end{equation}
Indeed, sum the $n$ second moments for $k=2$, and the at most
$n^2$ fourth moments for $k=4$, then apply Jensen's inequality.
For two parameter tuples at distance at most $\Delta$, the same
argument, using \eqref{eq:truncated-word-increment-lipschitz}, gives
\[
 \E N_k\bigl((\widehat Q_\theta-\E\widehat Q_\theta)
             -(\widehat Q_{\theta'}-\E\widehat Q_{\theta'})\bigr)
 \le C\sqrt\Delta.
\]
The random norms in this last display concentrate around their
means at Gaussian scale $C\sqrt{\Delta/n}$, again by
\eqref{eq:truncated-word-increment-lipschitz}.  The centered
expectations are deterministic and do not change Lipschitz
constants.

To control every parameter tuple, approximate it successively
on finer meshes.  For each $\ell\ge0$, choose a sup-norm mesh
$\mathcal N_\ell$ of spacing $2^{-\ell}$ in the parameter box.  It may be chosen with
\[
 |\mathcal N_\ell|\le\exp(C_L(\ell+1)n).
\]
For a tuple $\theta$, choose nearest mesh points
$\theta_\ell\in\mathcal N_\ell$.  Then
$\norm{\theta_\ell-\theta_{\ell-1}}_\infty\le C2^{-\ell}$.
For all pairs of mesh points at this distance, concentration and
a union bound show that their centered increments have both
$N_2$ and $N_4$ norm at most
\[
 C'\sqrt{\ell+1}\,2^{-\ell/2},
\]
except with probability $Ce^{-c(\ell+1)n}$, provided $C'$ is
chosen sufficiently large.  Specifically, the squared ratio
of this threshold to the Gaussian scale is a constant multiple
of $(\ell+1)n$, and its constant can be chosen to dominate the
logarithm of the number of mesh pairs.  The initial mesh is
bounded in the same way using
\eqref{eq:centered-word-norm-expectations}.  Summing the failure
probabilities over $\ell$, and summing the convergent increment
bounds along each chain, proves, with probability
$1-Ce^{-c_1n}$,
\begin{equation}\label{eq:simultaneous-centered-words}
 \sup_\theta N_k(\widehat Q_\theta-\E\widehat Q_\theta)\le C
 \quad(k=2,4).
\end{equation}
At each fixed $n$, continuity of the truncated word and its
expectation in $\theta$ justifies taking the limit of the mesh
chain.  There are only finitely many letter patterns of length
at most $L$, so their events may be intersected.

The centered words are now controlled on one event for every
parameter tuple.  Lemma~\ref{lem:word-bias} supplies the remaining
bias estimate,
$N_2(\E\widehat Q_\theta-\mathsf P(Q_\theta))\le C/\sqrt n$,
uniformly in $\theta$.  Sign symmetry makes the off-diagonal of
$\E\widehat Q_\theta$ identically zero.  Combining these facts
with \eqref{eq:simultaneous-centered-words} proves the desired
diagonal and off-diagonal assertions for the truncated words.
Their operator norms are already bounded by construction.
On $\norm W\le R$ and at parameters with $S_A(1;W)\succeq cI$,
the truncations equal the exact words, proving the GOE version.

The GOE estimate transfers to the zero-diagonal model through
$J=W-D_{\diag W}$.  This last step uses a Frobenius bound on the
removed diagonal.  Since
\[
 \norm{D_{\diag W}}_{\Frob}^{2}
   \ \stackrel{\mathrm{law}}=\ \frac{2j}{n}\chi_n^2,
\]
it is bounded by a fixed constant except with exponentially
small probability.  The norm tail for $W$, and a sufficiently
large choice of $R$, ensure also that both $W$ and $J$ lie in the
operator-norm ball except with exponentially small probability.
For every parameter tuple, \eqref{eq:truncated-word-lipschitz}
then bounds the Frobenius norm of the difference of the truncated
word evaluated at $W$ and at $J$ by a fixed constant.  Both
$N_2$ and $N_4$ of this difference are bounded by its Frobenius
norm.  The prediction is unchanged by this substitution.  At
parameters with $S_A(1;J)\succeq cI$, the truncated factors
evaluated at $J$ equal the exact factors.  This proves all of
\eqref{eq:word-estimates}.  The polynomial assertion was proved
in Lemma~\ref{lem:word-bias}.
\end{proof}

\subsection{Bilinear forms and regularized determinants}

We next extract two forms of the preceding estimates.  The first
controls bilinear forms for a prescribed finite family of
deterministic vectors.  The second controls a regularized
determinant simultaneously over its diagonal parameters and
bounded finite-rank perturbations.  Their quantifiers differ,
and we keep the two statements separate.

\begin{lemma}[Bilinear deterministic equivalents]\label{lem:bilinear}
Fix an integer $m\ge1$, a vector bound $M_1<\infty$, a tolerance
$\varepsilon>0$, and $R>2\sqrt j$.  For every deterministic
$0\le A\le a_+I$ and deterministic vectors
$v_1,\ldots,v_m\in\R^n$ with $\norm{v_r}\le M_1$, there is an
event of probability at least $1-Ce^{-c_3n}$ on which
\eqref{eq:deterministic-path-stability} holds and, with
\[
 G_A=A^{1/2}S_A(1;W)^{-1}A^{1/2}=K_A(W)A,
\]
one has, for every $r,s$,
\begin{align}
 |v_r^T(G_A-A)v_s|&\le\varepsilon,
 \label{eq:bilinear-G}\\
 |v_r^T(G_AW-B_AA)v_s|&\le\varepsilon,
 \label{eq:bilinear-GW}\\
 |v_r^T(WG_AW-B_AI-B_A^2A)v_s|&\le\varepsilon.
 \label{eq:bilinear-WGW}
\end{align}
Here $C,c_3,n_0$ may depend on $j,a_+,m,M_1,\varepsilon,R$ but
are uniform in $A$, the vectors, and $n\ge n_0$.  The assertion
is for each deterministic choice of these parameters; it does
not claim simultaneity over all vector choices.
\end{lemma}

\begin{proof}
For each of the specified bilinear forms, we calculate its
expectation and then apply concentration.  Use the truncations
above with a cutoff agreeing with the inverse on the path-stability
event.  Sign conjugation makes the off-diagonal expectations of
each relevant matrix vanish.
Equation~\eqref{eq:K-and-WK-bias} at $z=1$ gives
\[
 \E\widehat K_{ii}=1+O(n^{-1}),\qquad
 \E(\widehat W\widehat K)_{ii}=B_A+O(n^{-1}).
\]
The identities, valid on the event of agreement,
\[
 G_AW=K_A(W)(I+B_AA)-I,
 \qquad
 WG_AW=WK_A(W)(I+B_AA)-W
\]
therefore give the three diagonal expectation predictions
$A$, $B_AA$, and $B_AI+B_A^2A$, with entrywise errors
$O(n^{-1})$.  Their use with truncated factors changes
expectations only exponentially little, by the same argument
as in Lemma~\ref{lem:word-bias}.  In particular, their errors
in any of the specified bilinear forms are $O(n^{-1})$:
for a diagonal error matrix $E$, one has
$|v_r^TEv_s|\le\norm E\norm{v_r}\norm{v_s}$.

The expectation calculation has identified the three required
predictions.  To control fluctuations, each truncated bilinear form
is $C M_1^2$-Lipschitz in $W$ for Frobenius norm, by
\eqref{eq:truncated-word-lipschitz}.  Gaussian
concentration gives failure probability $Ce^{-c_3n}$ at any
fixed positive tolerance.  Intersect these events for the
$3m^2$ forms and with Lemma~\ref{lem:matrix-path-stability}.
For large $n$ the expectation errors are less than half the
specified tolerance, proving the result.
\end{proof}

\begin{lemma}[A simultaneous regularized determinant bound]
\label{lem:logdet}
Fix $\varepsilon>0$, an integer $r_0\ge0$, and $M_2<\infty$.
There is $\gamma_0>0$, depending only on $j,\varepsilon$, such
that for every fixed $0<\gamma\le\gamma_0$ and all sufficiently
large $n$, an event of probability at least $1-Ce^{-c_4n}$ has
the following property.  Simultaneously for all diagonal
$0\le V\le I$ and all real matrices $E$ with
$\rank E\le r_0$ and $\norm E\le M_2$, put
\[
 b=\frac1n\Tr V,\qquad B=jb,\qquad
 L=I+(BI-W)V+E.
\]
Then
\begin{equation}\label{eq:logdet-upper-bound}
 \frac1n\log\det(L^TL+\gamma I)\le jb^2+\varepsilon.
\end{equation}
The constants $C,c_4$ and the lower threshold on $n$ may depend
on $j,\varepsilon,\gamma,r_0,M_2$.  In particular, $V$ and $E$
may depend on $W$.
\end{lemma}

\begin{proof}
We first compute the mean log determinant with $E=0$ and a fixed
deterministic $V$.  Regularization will then provide a global
Lipschitz bound, allowing a finite mesh in $V$ and finally the
finite-rank perturbation.  Set
\[
 L_z=I+z^2BV-zWV,
 \qquad K(z)=(I-zVW+z^2BV)^{-1}.
\]
Since $V\le I$, we use the proof of
Lemma~\ref{lem:matrix-path-stability} with the diagonal upper bound
equal to $1$; its stability and inverse-norm constants here depend
only on $j$.
On the path-stability event from
Lemma~\ref{lem:matrix-path-stability} with $A=V$, the determinant
of $L_z$ equals the positive determinant of $S_V(z;W)$ by
$\det(I-UV)=\det(I-VU)$, and both
$L_z^{-1}$ and $K(z)$ have uniformly bounded operator norm.
Indeed $L_z=K(z)^{-T}$, and
\eqref{eq:K-symmetric-representation} applies.  Differentiating
the normalized log determinant on this event gives
\begin{align}
 \frac{\dd}{\dd z}\frac1n\log\det L_z
 &=\frac1n\Tr\bigl[L_z^{-1}(2zBV-WV)\bigr]\notag\\
 &=2zB\frac1n\Tr(K(z)V)
      -\frac1n\Tr(WK(z)V).
 \label{eq:logdet-derivative}
\end{align}
Here $VL_z^{-1}=K(z)V$, which follows by multiplying
$(I-zVW+z^2BV)V=V L_z$, and cyclicity of trace was used.

The derivative is expressed in the two matrices whose diagonal
expectations were computed in \eqref{eq:K-and-WK-bias}.
Let $\mathcal E_V$ be the path-stability event intersected with
$\norm W\le R$, for a fixed $R>2\sqrt j$.  This one event applies
to the entire interval of $z$.  Equations~\eqref{eq:K-and-WK-bias}, and
replacement by the truncated factors off $\mathcal E_V$, show
uniformly in $z,V$ that the expectation of the right-hand side
of \eqref{eq:logdet-derivative}, multiplied by
$\one_{\mathcal E_V}$, is
\[
 2zBb-zBb+O(n^{-1})=zjb^2+O(n^{-1}).
\]
All differentiated quantities are bounded on $\mathcal E_V$,
so integration and expectation may be interchanged.  Since
$L_0=I$, it follows that
\begin{equation}\label{eq:unregularized-logdet-mean}
 \E\left[\one_{\mathcal E_V}\frac2n\log\det L_1\right]
 =jb^2+O(n^{-1}).
\end{equation}
On $\mathcal E_V$, the singular values of $L_1$ are bounded
below by a positive constant.  Therefore
\[
 0\le\frac1n\log\det(L_1^TL_1+\gamma I)
       -\frac2n\log\det L_1
 \le C\gamma.
\]
Off $\mathcal E_V$, the absolute value of the regularized
normalized log determinant is bounded by
\[
 |\log\gamma|+C+2\log(1+\norm W),
\]
for $0<\gamma\le1$.  The norm tails of $W$ and
\eqref{eq:deterministic-path-stability} make its expected
contribution exponentially small, for each fixed $\gamma$.
Consequently, uniformly in deterministic $V$,
\begin{equation}\label{eq:regularized-logdet-mean}
 \E\left[\frac1n\log\det(L_1^TL_1+\gamma I)\right]
 \le jb^2+C\gamma+O(n^{-1})+O_\gamma(e^{-c n}).
\end{equation}

The mean bound is uniform in deterministic $V$.  To obtain a
single event for all $V$, we use the concentration supplied by
regularization.  Define on all real matrices
\[
 F_\gamma(M)=\frac1n\log\det(M^TM+\gamma I).
\]
Its Frobenius gradient is
$2n^{-1}M(M^TM+\gamma I)^{-1}$.  Each singular value of this
gradient is at most $1/(n\sqrt\gamma)$, and hence
\begin{equation}\label{eq:regularized-logdet-lipschitz}
 \Lip_{\Frob}(F_\gamma)\le\frac1{\sqrt{n\gamma}}.
\end{equation}
For fixed $V$, the map $W\mapsto L_1$ has Frobenius Lipschitz
constant at most one.  Gaussian concentration therefore gives
\begin{equation}\label{eq:regularized-logdet-concentration}
 \Prob\{|F_\gamma(L_1)-\E F_\gamma(L_1)|>u\}
 \le2\exp\left(-\frac{\gamma n^2u^2}{4j}\right).
\end{equation}

Choose $\gamma_0\le1$ small enough that $C\gamma_0$ in
\eqref{eq:regularized-logdet-mean} is at most
$\varepsilon/8$.  For fixed $\gamma\le\gamma_0$ and sufficiently
large $n$, the expectation bound and
\eqref{eq:regularized-logdet-concentration} bound all points
of any fixed sup-norm mesh of the cube $[0,1]^n$ by
$jb^2+\varepsilon/2$, with failure $\exp(-c_\gamma n^2)$.
Indeed a mesh of spacing $\Delta>0$ has at most
$(C/\Delta)^n$ points, whose logarithm is only of order $n$.
On $\norm W\le R$, if $\norm{V-V'}\le\Delta$, then
\[
 |b-b'|\le\Delta,\qquad
 \norm{(I+(BI-W)V)-(I+(B'I-W)V')}_{\Frob}
 \le\sqrt n(2j+R)\Delta.
\]
By \eqref{eq:regularized-logdet-lipschitz}, the corresponding
$F_\gamma$ values differ by at most
$(2j+R)\Delta/\sqrt\gamma$, whereas $|jb^2-j(b')^2|\le2j\Delta$.
Choose $\Delta$, depending only on
$j,R,\gamma,\varepsilon$, so that these errors sum to at most
$\varepsilon/4$.  Including the exponentially likely norm
event proves, simultaneously for all $V$,
\[
 F_\gamma(I+(BI-W)V)\le jb^2+3\varepsilon/4.
\]

We have obtained the simultaneous bound for $E=0$ with room for
the perturbation.  The same Lipschitz estimate
\eqref{eq:regularized-logdet-lipschitz} gives, for every allowed $E$,
\[
 |F_\gamma(M+E)-F_\gamma(M)|
 \le\frac{\norm E_{\Frob}}{\sqrt{n\gamma}}
 \le M_2\sqrt{\frac{r_0}{n\gamma}}.
\]
This is at most $\varepsilon/4$ for large $n$, uniformly even
when $E$ is chosen from the observed matrix.  It proves
\eqref{eq:logdet-upper-bound} and the lemma.
\end{proof}
